\chapter{Scabies}

\section*{Definition}
Scabies is a contagious skin infestation caused by the mite Sarcoptes scabiei var. hominis. It causes intense itching, especially at night, and a characteristic rash. It spreads through close skin-to-skin contact and is common worldwide.

\section*{What Happens in Your Body}
The female mite burrows into the stratum corneum, creating tunnels where she lays eggs. The burrows trigger an immune response (type IV hypersensitivity) to the mites, eggs, and feces. Sensitization takes 2-6 weeks after initial infestation, so symptoms are delayed. The intense itch is primarily driven by the immune response, not the mites themselves.

\section*{Causes}
\begin{itemize}
\item Close skin-to-skin contact with an infested person
\item Sexual contact
\item Sharing clothing, bedding, or towels
\item Crowded living conditions (nursing homes, prisons, dormitories)
\item Immunosuppression (HIV, corticosteroid use) leading to crusted scabies
\item Poverty and poor hygiene
\end{itemize}

\section*{When to See a Doctor}
\begin{itemize}
\item Intense itching, worse at night
\item Thin, grayish-white, S-shaped burrows (most common between fingers, wrists, elbows, axillae, waist, genitals)
\item Papules, vesicles, and excoriations
\item Red, raised bumps (nodules) in the axillae and groin
\item In infants: rash may involve the face, scalp, palms, and soles
\item Crusted scabies: thick, crusted plaques with minimal itching (in immunocompromised)
\item Multiple household members affected
\end{itemize}

\section*{Related Symptoms}
\begin{itemize}
\item Intense itching (pruritus)
\item Eczema (can be secondarily eczematized)
\item Contact dermatitis
\item Folliculitis
\item Impetigo (secondary bacterial infection)
\end{itemize}

\section*{Self-Care / Home Management}
\begin{itemize}
\item Wash all clothing, bedding, and towels in hot water and dry on high heat
\item Items that cannot be washed should be sealed in a plastic bag for 72 hours
\item Notify close contacts and household members for treatment
\item Do not use OTC scabicides (permethrin is prescription)
\item Avoid scratching to prevent secondary infection
\end{itemize}

\section*{How Your Doctor Will Evaluate You}
\begin{itemize}
\item Clinical diagnosis based on history of itching and characteristic burrows
\item Microscopic examination of skin scraping from a burrow (mites, eggs, feces visible)
\item Dermoscopy to visualize the mite at the end of a burrow
\end{itemize}

\section*{Treatment Options}
\begin{itemize}
\item Permethrin 5 percent cream (first-line): apply from neck to toes, wash off after 8-14 hours
\item Ivermectin (oral): for resistant cases, crusted scabies, or institutional outbreaks
\item Sulfur ointment for infants and pregnant women
\item Lindane (second-line, due to neurotoxicity risk)
\item Treat all household members and close contacts simultaneously
\item Antihistamines for itching
\item Topical corticosteroids for persistent nodules
\end{itemize}

\section*{References}
\begin{thebibliography}{99}
\bibitem{scabies:who-scabies} Engelman D, et al. "The 2023 International Alliance for the Control of Scabies classification." \emph{Lancet Infect Dis}. 2023;23(8):e298-e308.
\bibitem{scabies:uptodate-scabies} Chosidow O, et al. Scabies: epidemiology, clinical features, and diagnosis. In: UpToDate, Post TW (Ed), UpToDate, Waltham, MA; 2025.
\bibitem{scabies:nejm-scabies} Chosidow O. "Scabies." \emph{N Engl J Med}. 2023;388(13):1212-1222.
\bibitem{scabies:bmj-scabies} Engelman D, et al. "Scabies." \emph{BMJ}. 2023;382:e075549.
\bibitem{scabies:cochrane-scabies} Rosumeck S, et al. "Ivermectin versus permethrin for scabies." \emph{Cochrane Database Syst Rev}. 2023;(6):CD012995.
\end{thebibliography}
